\documentclass[11pt]{article}
\usepackage[T1]{fontenc}
\usepackage{lmodern,amsmath,amssymb,amsthm,microtype,booktabs}
\usepackage[letterpaper,margin=1in]{geometry}
\usepackage[hidelinks]{hyperref}
\usepackage{fancyhdr}
\hypersetup{pdftitle={Prime outputs and local counts for two shifted Mills minima},pdfauthor={Research note for Keio}}
\pagestyle{fancy}\fancyhf{}\lhead{Prime outputs and local counts}\rhead{Shifted Mills minima}\cfoot{\thepage}\setlength{\headheight}{14pt}
\newtheorem{theorem}{Theorem}
\newcommand{\fl}[1]{\lfloor #1\rfloor}
\newcommand{\RR}{\operatorname{Re}}
\newcommand{\li}{\operatorname{li}}
\title{Prime outputs and local counts for two shifted Mills minima}
\author{Research note for Keio}\date{November 3, 2026}
\begin{document}\maketitle
\begin{abstract}
Two fixed exponent sequences give simple recurrences and certified large primes. A three-term rational seed produces seven consecutive pairs ending in 400- and 401-digit primes. A second seed reaches 802 and 803 digits. For the separate least bases, exact eight-level computations give unconditional lower bounds. Under RH they identify the initial outputs of the infinite minima and the first 60 digits of two primes of more than 263 million digits. We also show how local counts at the heights supplied by just one minimum detect the supremum of the real parts of zeta zeros. This last result uses all ordinary primes in the indicated windows. The required square-root error estimate remains unproved.
\end{abstract}

\section{What the two least bases determine}
Write $E_k^e=3^k+e$ for $e\in\{-1,1\}$ and $k\ge1$. Let $A_e$ be the least real $A>1$ with all $p_k^e=\fl{A^{E_k^e}}$ prime. Existence follows from Saito's theorem \cite{saito}. Both transcendence statements are proved in the earlier corpus work and accompanying research note \cite{prior}. We distinguish (i) these infinite minima, (ii) minima subject to finitely many prime conditions, and (iii) explicitly chosen bases used for two finite sets of conditions at once.

For either infinite minimum, the exact relation $E_{k+1}^e=3E_k^e-2e$ gives a limit of rational expressions in prime outputs,
\[
 \frac{p_{k+1}^e}{(p_k^e)^3}=A_e^{-2e}+O_{A_e}(1/p_k^e).
\]
Indeed $p_k^e\le A_e^{E_k^e}<p_k^e+1$ and taking a floor changes a real number by less than one. Each limit is transcendental. The two numbers $A_-A_+$ and $A_-/A_+$ cannot both be algebraic, since their product is $A_-^2$. More generally, among $A_-^jA_+$ for integral $j$ at most one can be algebraic: the quotient of two exceptions would be a nonzero power of $A_-$. These deductions do not determine whether a particular quotient is irrational or whether the pair is algebraically independent. The Saito--Takeda independence theorem concerns other families \cite{st}.

There is an exact simplification in the direction of smaller exponents. Suppose $P=\fl{A^N}$ is known and $q=\fl{A^M}$ is known to be an odd prime with $0<M<N$. Then
\[
 \boxed{q=2\fl{P^{M/N}/2}+1.}
\]
The interval $[P^{M/N},(P+1)^{M/N})$ has length less than one. Its possible floors are the floor of its left endpoint and the next integer. Exactly one of those two is odd. The argument works for outputs on either schedule or both schedules sharing the same base. Thus, for example, the single 803-digit prime below, with its exponent and the list of lower exponents, recovers all thirteen earlier odd prime values without knowing the base. To evaluate forward from a certified interval $[L,U]$ for a minimum, the same rule applies whenever $U^N-L^N<1$ and the target prime exceeds two. Transcendence by itself gives no numerical enclosure of this precision.

\section{Producing large prime pairs}
Use one common trial base $A=\sqrt B$. Put $h_k=(3^k-1)/2$, $t_1=B$ and $t_{k+1}=Bt_k^3$. Since $h_{k+1}=3h_k+1$, the pair from the two original exponents is exactly
\[
       (\fl{t_k},\fl{Bt_k})=(\fl{B^{h_k}},\fl{B^{h_k+1}}).
\]
The following rational input is written with exact terminating decimals:
\[
 B=2.32167913260277151944-337679\,10^{-52}+978854\,10^{-403}.
\]
The first three pairs are $(2,5)$, $(29,67)$, $(56941,132199)$. All fourteen outputs at levels one through seven are primes.
\begin{center}
\begin{tabular}{crrrrrrr}
\toprule
Level &1&2&3&4&5&6&7\\
\midrule
Digits at $3^k-1$ &1&2&5&15&45&134&400\\
Digits at $3^k+1$ &1&2&6&15&45&134&401\\
\bottomrule
\end{tabular}
\end{center}
This seed was found by searching inside a finite intersection of prime floor intervals and then certifying the resulting integers. In a second search we found a rational $A$ with denominator $10^{820}$, beginning $2.3260749786289520208\ldots$, for which the seven pairs start $(5,29)$ and end with primes of 802 and 803 digits. At the very same $A$, the next two floors, at exponents 6560 and 6562, have proper factors 163 and 31 respectively. The accompanying certificate archive contains the full numerator, all exact integers, and elliptic-curve certificates checked with standard Python. Primality is unconditional; the two trial bases are finite constructions. The recurrence evaluates powers efficiently but does not make an arbitrary seed generate primes indefinitely.

\section{What can be identified for the actual minima?}
Let $L_{e,j}$ denote the least admissible base through level $j$. Exact finite search gives upper and lower endpoints for the last prime cell of each exhibited prefix. To justify minimality inductively, start at the preceding least base $L$. Find the first prime $q\ge\fl{L^{E_j}}$ and test all prior floors at $L'=\max(L,q^{1/E_j})$. If the tests pass, $L'$ is the new finite minimum: the skipped floors cannot yield a prime, and bases smaller than $L$ already failed an earlier test. All such tests pass here; this rule need not succeed for an arbitrary initial branch. Lower endpoints and upper endpoints are compared by crossing integer powers, so no numerical root decisions are needed.
\begin{center}
\begin{tabular}{lrr}
\toprule
 & negative shift & positive shift\\
\midrule
First three prime values & $2,17,10007$ & $2,5,131$\\
Digits of level-eight prime $q_e$ & 1,010 & 497\\
Exponent $E_8^e$ & 6,560 & 6,562\\
\bottomrule
\end{tabular}
\end{center}
Each separate last cell $[q_e^{1/E_8^e},(q_e+1)^{1/E_8^e})$ lies inside every prior cell of that prefix. Both lower endpoints are unconditional lower bounds for the respective infinite minima. A proof of continuation from both concrete cells would suffice for the numerical identifications below.

\paragraph{Continuation under RH.}
Schoenfeld's explicit estimate under RH is $|\vartheta(x)-x|<\sqrt{x}(\log x)^2/(8\pi)$ for $x\ge599$ \cite{schoenfeld}. Given $p\ge10^{18}$ and any successive exponent ratio $r\in[5/2,4]$, take $x=p^r$ and $y=x+p^{r-1}$. The sum of the two error bounds at $x$ and $y$ is at most $3p^{r/2}\log^2p$: use $y\le2x$, $\log y\le 5\log p$. Since $p^{r/2-1}\ge p^{1/4}>3\log^2p$ throughout this range, $\vartheta(y)>\vartheta(x)$. A prime $x<q\le y$ has $q+1<(p+1)^r$. Its entire next cell fits in the previous one. Nested closed subcells with strict next containments give an infinite continuation. Applied to either terminal cell, this proves under RH that its interval contains the infinite minimum. The certified finite primes themselves require no such assumption.

Write $E=3^{20}+e$ and $d=E_8^e$. The two logarithms $(E/d)\log_{10}q_e$ and $(E/d)\log_{10}(q_e+1)$ enclose $\log_{10}(A_e^E)$ under continuation. Each lower endpoint alone applies unconditionally. The resulting digit counts are
\begin{center}
\begin{tabular}{lrr}
\toprule
 & $p_{20}^{-}$ & $p_{20}^{+}$\\
\midrule
Unconditional digit lower bound & 536,469,125 & 263,660,230\\
Exact number under RH & 536,469,125 & 263,660,230\\
First 20 digits under RH & 22248781440454750739 & 58063115119761947324\\
\bottomrule
\end{tabular}
\end{center}
No expansion of either giant integer is needed for these bounds. For the full 60-digit leading integer $h$ supplied in the certificate, with digit count $d_{20}$, the verifier establishes by outward rational logarithm bounds that
\[
 h\,10^{d_{20}-60}<q_e^{E/d}<(q_e+1)^{E/d}
                 <(h+1)10^{d_{20}-60}.
\]
In particular the lower bound remains valid even if the terminal cell fails to extend.

\paragraph{A relation between the two bases.}
Under the same continuation condition, the ratio $R=\log A_-/\log A_+$ lies in an explicit interval beginning $2.0346986920834196325654\ldots$. Substituting the lower logarithm for the numerator and upper for the denominator gives a lower endpoint, and reversing the choices gives an upper endpoint. The archive verifies, using integer bounds on logarithms, that this enclosure is inside a Farey interval $(a/b,c/d)$ with $bc-ad=1$ and $b+d>2.51\cdot10^{249}$. Any reduced rational in that interval has denominator at least $b+d$, since $bu-av$ and $cv-du$ are positive integers for $u/v$ inside. In particular, under RH the equation $A_-^v=A_+^u$ is impossible for positive integers $v\le 10^{249}$. Larger rational relations and the unconditional order of the two bases are not decided by this computation.

\section{A test using the other primes near a single output}
Fix either sequence $p_k=p_k^e$. Take $X=p_k$, $Y=X/\log X$, and define
\[
 W(X)=\sup_{Y\le u\le X}
 \left|\pi(u)-\pi(Y)-\int_Y^u\frac{dt}{\log t}\right|.
\]
Here and below $X$ is sufficiently large. This statistic uses only the ordinary primes in $(Y,X]$ and no cumulative count below $Y$.
\begin{theorem}\label{th:obs}
RH holds if and only if $W(p_k)\le \sqrt{p_k}\log p_k$ for infinitely many $k$. Under RH this inequality holds for every sufficiently large $k$. If RH fails then $W(p_k)/(\sqrt{p_k}\log p_k)\to\infty$. The statements hold for either fixed minimum separately and for any other sequence $X_k\to\infty$.
\end{theorem}
Under RH the classical explicit bound for $|\pi(x)-\li(x)|$ is $\sqrt{x}\log x/(8\pi)$ for $x\ge2657$ \cite{schoenfeld,dlmf}. Applying it at $u$ and $Y$ proves the upper bound. For the converse we give the following local version of the zero-detection argument. Lower bounds for oscillations of $\psi(x)-x$ isolated from a fixed zeta zero, including localization and mean estimates, have a substantial history in work of Tur\'an and Pintz \cite{pintz84,pintz22}. We do not assert historical priority for localization. Centering the statistic at $Y$ is essential for using only the primes inside the window.

\paragraph{The local discrepancy exponent.}
In fact, if $\Theta=\sup\{\RR\rho:0<\RR\rho<1,\ \zeta(\rho)=0\}$, then
\[
 \lim_{X\to\infty}\frac{\log(2+W(X))}{\log X}=\Theta .
\]
We prove this for the anchored discrepancy; the result also proves the converse in Theorem~\ref{th:obs}. Let $\rho$ be a nonreal zero with $\beta=\RR\rho\ge1/2$, of multiplicity $m$. Write $D(s)=-\zeta'(s)/(s\zeta(s))-1/(s-1)$. For $\RR s>1$ this is the Mellin transform of $\Delta(x)=\psi(x)-x$ on $x\ge1$ in the convention $\int f(v)v^{-s}dv/v$. Put
\[
 \xi(s)=\tfrac12s(s-1)\pi^{-s/2}\Gamma(s/2)\zeta(s),\qquad
 G(s)=\frac{s(s-1/2)\xi(s)}{(s-\rho)^m}.
\]

The inverse Mellin transform $h$ of $G$ satisfies
$|h(v)|\ll_\rho\exp(-2\max(v^2,v^{-2}))$. Indeed the inverse of $\pi^{-s/2}\Gamma(s/2)\zeta(s)$ on $\RR s>1$ is
$2\sum_{n\ge1}\exp(-\pi n^2/v^2)$, by termwise integration. Apply $(v\partial_v)(v\partial_v-1)/2$ to obtain the kernel $g$ for $\xi$. The Gaussian summation formula gives $g(v)=vg(1/v)$, because the differential operator annihilates the constant and linear terms. The series near zero and this symmetry give Gaussian decay with exponential constant $\pi$ at both ends, up to polynomial factors, also for all logarithmic derivatives. The numerator factors $s(s-1/2)$ act by such derivatives. To divide by $s-\rho$, replace the current kernel $k$ by
\[
 v^\rho\int_0^v k(w)w^{-\rho}\,\frac{dw}{w}
 =-v^\rho\int_v^\infty k(w)w^{-\rho}\,\frac{dw}{w}.
\]
At each of the $m$ steps the old transform vanishes at $\rho$. Use the first expression for small $v$ and the second for large $v$. Integrating by parts verifies the transform, and the two Gaussian tail bounds persist up to polynomial factors. Since $2<\pi$, this proves the stated bound on $h$.

Take $Z=X/\sqrt{\log X}$. In $-1\le\RR s\le2$ the product $GD$ has only the simple pole at $\rho$, with nonzero residue $c_\rho$. The nontrivial zero factors cancel through $\xi$; $D$ is regular at one and $G(0)=0$. On the displaced contours, cancellation of $\zeta$ in the logarithmic derivative and Stirling's formula give exponential vertical decay. Thus convolution on $\RR s=2$ followed by a shift to $-1$ gives
\[
 \int_1^\infty \Delta(u)h(Z/u)\frac{du}{u}=c_\rho Z^\rho+O(Z^{-1}).
\]
We may replace $\psi(u)$ here by $\vartheta(u)$ to within $o(\sqrt Z)$: by the ordinary prime number theorem, $\psi(u)-\vartheta(u)=\vartheta(\sqrt u)+O(u^{1/3}\log^2(2u))=\sqrt u+o(\sqrt u)$. The transform vanishes at $1/2$, so the full $\sqrt u$ contribution is zero (truncation at $u=1$ is exponentially small). For the remainder substitute $v=Z/u$ and apply dominated convergence against $v^{-1/2}|h(v)|dv/v$.

Now $Z/Y=X/Z=\sqrt{\log X}$. With $\vartheta(u)-u=O(u)$, discarding $u<Y$ and $u>X$ costs $o(1)$. For example the upper error is bounded by $CZ\int_{\sqrt{\log X}}^\infty e^{-2t^2}dt$ with $t=u/Z$. The lower tail follows from the same estimate with $t=Z/u$. Since $G(0)=0$, subtracting the unknown endpoint value $\vartheta(Y)-Y$ from the remaining integrand also costs $o(1)$. The kernel has bounded $L^1(dv/v)$ norm. It follows for every sufficiently large $X$ that
\[
 W_{\vartheta}(X):=\sup_{Y\le u\le X}|\vartheta(u)-\vartheta(Y)-(u-Y)|
       \gg_\rho (X/\sqrt{\log X})^\beta .
\]
Local partial summation gives $W_{\vartheta}(X)\le2(\log X)W(X)$ and $W(X)\le W_{\vartheta}(X)/\log Y$. For the matching upper exponent, the classical truncated explicit formula at half-integers gives $\psi(u)=u+O(u^\Theta\log^2(2u))$, using a truncation height comparable to $u$ and $\sum_{|\Im\rho|\le 2u}|\rho|^{-1}\ll\log^2(2u)$ \cite{mv}. Monotonicity then brackets arbitrary real inputs; for $\Theta=1$ the trivial $O(u)$ bound suffices. Passing to $\vartheta$ does not increase the exponent. The lower bound holds for each zero with $\beta\ge1/2$, including a critical-line zero, and proves the claimed limit. If RH fails, choose $\beta>1/2$ using the functional equation; its lower bound diverges after division by the threshold in Theorem~\ref{th:obs}.

If one wants the observable base parameters explicitly, at level $E_k^e$ every integer in $(Y,p_k]$ can be tested as a raw floor using bases in $[\min(Y^{1/E_k^e},A_e),A_e]$. The left endpoint tends to $A_e$, and the length is asymptotic to $A_e\log(E_k^e)/E_k^e$. Here the base can change at every query. The single digit $p_k$ determines the window; knowing that $A_e$ is transcendental gives no upper bound for its interior prime discrepancy.

\section{A limit on inference from local observations}
For comparison, permit a subset $P$ of the genuine primes as the admissible prime set, writing $\pi_P$ for its count. There exist fixed open neighborhoods of both actual minima such that the indicators of $P$ and of the primes agree on every raw floor from those neighborhoods at either schedule. Even the cumulative counts at all those heights $y$ can satisfy $\pi_P(y)-\pi(y)=O(y^{9/20})$, while
\[
       |\pi_P(x)-\li(x)|\gg x^{3/4}
\]
along an unbounded sequence, with $\pi_P(x)\sim x/\log x$.

To see this, consider the two points $\log\log A_e$ modulo $\log3$. There is a complementary arc of length at least $(\log3)/2>\log(5/3)$. Take the neighborhoods small enough that after projection a slightly enlarged arc of length $\log(5/3)$ remains vacant. Uniformly on their compact closures,
$\log\log\fl{A^{3^k+e}}=k\log3+\log\log A+o(1)$.
Consequently there are unobserved intervals $[T_j,T_j^{5/3}]$ with $T_{j+1}=T_j^3$. Discard finitely many initial intervals. In $[T_j,2T_j]$ the ordinary prime number theorem allows deletion of either $n_j=\fl{T_j^{3/4}}$ or $2n_j$ primes. Choose the number that leaves the count at $2T_j$ at distance at least $n_j/2$ from $\li(2T_j)$. Only $O(x^{3/4})$ primes have been deleted through $x$, preserving the prime number theorem. At any observed height $y$ after block $j$ and before the next deletion, $y>T_j^{5/3}$, so the total lost count is $O(T_j^{3/4})=O(y^{9/20})$. Both global minima are unchanged because they remain admissible and deletion cannot introduce a smaller base.

This model agrees even on entire candidate neighborhoods. Its prime set differs from the ordinary primes and its counting function is not subject to the zeta-function proof in Section 4. For the ordinary primes, computing or proving the square-root bound inside infinitely many of the stated windows remains an RH-level step.

\appendix
\section{Reproducing the prime outputs}
The following exact Python calculation prints the complete integers for the three-term seed. The numerator and denominator of $B$ have a common decimal scale of $10^{403}$. Elliptic-curve and Lucas proof verification is separate from the floor evaluation.
\begin{verbatim}
b = (232167913260277151944 * 10**383
     - 337679 * 10**351 + 978854)
d = 10**403
for k in range(1, 8):
    n = (3**k - 1)//2
    print(k, b**n // d**n, b**(n+1) // d**(n+1))
\end{verbatim}
The archive's \texttt{prime\_outputs.txt} lists the full pair values from both trial bases and all original finite prefix outputs. Its standard-library scripts independently verify primality, finite minimality, the two composite successor floors for the second trial base, and the outward bounds on the leading digits and logarithm ratio. For instance, the 60 leading digits conditional on continuation at level twenty are as follows, in blocks of ten:
\begin{center}\small
\begin{tabular}{l@{\quad}l}
$-$ & \texttt{2224878144 0454750739 5301073243 9960956130 0507646826 4118478769}\\
$+$ & \texttt{5806311511 9761947324 6743688848 0153219132 4720382463 4113594065}
\end{tabular}
\end{center}
\begin{thebibliography}{9}
\bibitem{prior} Keio, \emph{Ternary trace dynamics and the transcendence of shifted Mills constants}, corpus manuscript 064 (2026); see also \emph{Two fixed Mills minima and the zeros of Euler products}, research note of November 3, 2026.
\bibitem{saito} K. Saito, Transcendency of variants of Mills' constant, \emph{Ramanujan J.} 70 (2026), 69. \href{https://arxiv.org/html/2508.16068v3}{arXiv:2508.16068v3}.
\bibitem{st} K. Saito and W. Takeda, Topological properties and algebraic independence of sets of prime-representing constants, \emph{Mathematika} 68 (2022), 429--453. \href{https://doi.org/10.1112/mtk.12142}{doi:10.1112/mtk.12142}.
\bibitem{schoenfeld} L. Schoenfeld, Sharper bounds for the Chebyshev functions $\theta(x)$ and $\psi(x)$. II, \emph{Math. Comp.} 30 (1976), 337--360. \href{https://doi.org/10.1090/S0025-5718-1976-0457374-X}{doi:10.1090/S0025-5718-1976-0457374-X}.
\bibitem{dlmf} NIST Digital Library of Mathematical Functions, \href{https://dlmf.nist.gov/27.12.E7}{equation 27.12.7}.
\bibitem{mv} H. L. Montgomery and R. C. Vaughan, \emph{Multiplicative Number Theory I: Classical Theory}, Cambridge University Press (2007), chapters 12--15.
\bibitem{pintz84} J. Pintz, \emph{On the remainder term of the prime number formula and the zeros of Riemann's zeta-function}, Lecture Notes in Mathematics 1068 (1984), Theorem 2.
\bibitem{pintz22} J. Pintz, On the mean value of arithmetic error terms, \emph{Math. Pannon.} 28 (2022), 58--64. \href{https://doi.org/10.1556/314.2022.00007}{doi:10.1556/314.2022.00007}.
\end{thebibliography}
\end{document}
